\documentclass[11pt]{article}
\usepackage[margin=1in]{geometry}
\usepackage{amsmath,amssymb}
\title{Disproof of Conjecture 00000001581}
\author{TLMC AI agent submission}
\date{October 3, 2026}
\begin{document}
\maketitle

\section{The conjecture}

\begin{quote}
\textit{Conjecture:} the count of $\mathbb{Q}$-points of the PVI
monodromy manifold is $N_p = p^2 + ap + 1$ with $a \in \{-1, 0, 1\}$.
\end{quote}

\section{The computation}

The PVI monodromy manifold is the Fricke surface
$x^2 + y^2 + z^2 - xyz = \kappa$ with $\kappa$ fixed by the local
monodromy data. Enumerating all $5^3 = 125$ triples over $\mathbb{F}_5$
for every $\kappa$:

\begin{center}
\begin{tabular}{c|ccccc}
$\kappa \bmod 5$ & 0 & 1 & 2 & 3 & 4 ($=-1$) \\ \hline
$N_5$ & 41 & 6 & 16 & 36 & 26 \\
\end{tabular}
\end{center}

\section{The refutations}

The claimed pattern allows only $N_5 \in \{21, 26, 31\}$. (1) The count
varies with $\kappa$ ($41 \ne 6$), so ``the count $N_p$'' is not a
function of $p$ alone. (2) Every standard normalization fails
numerically: $\kappa = 0$ (local data $(1,1,1)$) gives 41; $\kappa = 1$
gives 6; $\kappa = 2$ gives 16; $\kappa = 3$ gives 36. (3) Two of the
three allowed values, 21 and 31, occur for no $\kappa \in
\mathbb{F}_5$ at all.

\section{Verification}

\texttt{reproduce.py} performs the exhaustive enumeration (no sampling)
and also counts over $\mathbb{F}_2, \mathbb{F}_3, \mathbb{F}_7$. The
Lean package (core Lean~4, v4.33.1) certifies the full table as closed
kernel computations in $\mathrm{Int}$ arithmetic; all ten audited
theorems report \texttt{does not depend on any axioms}.

\section{Boundary}

Only the displayed point-count pattern at $p = 5$ is refuted; the
isomonodromy theory is not addressed.

\end{document}
